The Python side is responsible for generating the test cases before they are executed on the FPGA verification platform. Instead of embedding predefined test scenarios directly into the hardware test logic, the verification inputs are constructed on the host computer and then converted into a ROM initialization file. This makes the test cases easier to modify, reproduce, and extend.

In this verification flow, the test-generation scripts define the packet-level information used by the DUT, including the source node, destination node, packet length, start time, message type, and flit contents. The generated packets are first organized as structured test cases and are then encoded into a ROM initialization format. This ROM data is used to initialize the on-chip memory before FPGA execution, allowing the FPGA-side controller to read predefined packet entries during the test.

This method separates test construction from the hardware verification logic. New packet sequences can be generated by modifying the Python case definitions, while the FPGA-side controller, generator, and checker can remain unchanged. Therefore, the same hardware verification platform can be reused for different test scenarios, such as single-flit packets, concurrent packet injection, multi-flit packets, and larger repeated test sets.

Overall, Python-based test case generation provides a flexible and reproducible way to prepare verification inputs. It allows complex packet scenarios to be described at a higher level, encoded automatically, and executed by the FPGA-side verification platform without changing the DUT or the on-chip test logic.
